\hypertarget{buxe1o-cuxe1o-thu-houx1ea1ch-ux1ee9ng-dux1ee5ng-agentic-ai-vuxe0-trux1ea3i-nghiux1ec7m-thux1ef1c-chiux1ebfn-hackathon}{%
\section{Báo cáo Thu hoạch: Ứng dụng Agentic AI và Trải nghiệm Thực
chiến
Hackathon}}

\hypertarget{mux1ee5c-tiuxeau-sux1ef1-kiux1ec7n}{%
\subsection{1. Mục tiêu Sự
kiện}}

\begin{itemize}
\tightlist
\item
  Trải nghiệm hành trình dấn thân, đối mặt với áp lực thời gian và rèn
  luyện tư duy phát triển sản phẩm tinh gọn (MVP) thông qua sức nóng của
  cuộc thi Agentic AI Build Week Hackathon.
\item
  Khám phá sức mạnh đột phá của Agentic AI trong việc giải quyết các bài
  toán vận hành thực tế của doanh nghiệp, trải dài từ kiểm soát đám
  đông, đặt món ăn đa kênh đến phân tích chiến lược rủi ro và hỗ trợ
  thiết kế kiến trúc hệ thống đám mây.
\end{itemize}

\hypertarget{diux1ec5n-giux1ea3-khuxe1ch-mux1eddi-cuxe1c-ux111ux1ed9i-thi}{%
\subsection{2. Diễn giả Khách mời (Các Đội
thi)}}

\begin{itemize}
\tightlist
\item
  \textbf{Team 3KA} (Huỳnh An Khương, Nguyễn Quốc Huy, Ngô Quang Khôi,
  Hoàng Lê Thành Đức, Đặng Nguyễn Phước Lộc, Đặng Trường Hưng) - Dự án
  S.H.E.P.H.E.R.D (Giám sát và dự báo dòng người).
\item
  \textbf{One Team} (Anh Duy, Tran Dong, Doan Trung, Minh Viet, Anshul
  Roy) - Dự án KFC Bot Agent (Trợ lý đặt món đa kênh).
\item
  \textbf{Plan V} (Pham Tien Thuan Phat, Huynh Hoang Long, Le Minh
  Nghia, Tran Dai Vi, Nguyen An) - Dự án SA Professional AI Native App
  (Tác tử hỗ trợ kiến trúc sư đám mây).
\item
  \textbf{Team Signal Scout} (Le Tan Luc, Do Hoang Hieu, Trieu Quoc Hao,
  Nguyen Van Duy Khiem, Nguyen Cong Minh, Nguyen Tran Minh Quan) - Dự án
  Signal Scout (Hệ thống phân tích rủi ro doanh nghiệp).
\end{itemize}

\hypertarget{nhux1eefng-ux111iux1ec3m-nhux1ea5n-quan-trux1ecdng-key-highlights}{%
\subsection{3. Những Điểm Nhấn Quan Trọng (Key
Highlights)}}

\hypertarget{sux1ef1-bux1ee9t-phuxe1-cux1ee7a-agentic-ai-khi-ai-biux1ebft-tux1ef1-lux1eadp-kux1ebf-houx1ea1ch-vuxe0-huxe0nh-ux111ux1ed9ng}{%
\subsubsection{Sự bứt phá của Agentic AI: Khi AI biết tự lập kế hoạch và
hành
động}}

Agentic AI đánh dấu một bước chuyển dịch mang tính cách mạng khi AI
không còn chỉ phản hồi thụ động. Thay vào đó, AI đã có khả năng tự động
lập kế hoạch, gọi các công cụ nghiệp vụ (tools/APIs) để giải quyết những
bài toán phức tạp và xác minh kết quả một cách độc lập.

Điều này được minh chứng rất rõ nét qua dự án KFC Bot Agent --- hệ thống
giúp khách hàng đặt món trực tiếp trên Zalo/Messenger mà không cần
chuyển app hay đăng ký rườm rà. Hoặc ấn tượng hơn là ứng dụng SA
Professional có thể tự động bóc tách yêu cầu nghiệp vụ (BRD), vẽ sơ đồ
kiến trúc trên Draw.io và ước tính chi phí hạ tầng AWS chỉ trong vài
phút.

\hypertarget{huxe0nh-truxecnh-24h-hackathon-ux111ux1ea7y-hux1ed7n-loux1ea1n-nhux1b0ng-vuxf4-giuxe1}{%
\subsubsection{Hành trình 24h Hackathon: Đầy hỗn loạn nhưng vô
giá}}

Cuộc thi Hackathon thực sự là một cuộc ``hành xác'' ngọt ngào và đầy cảm
xúc, đưa các đội thi qua đủ mọi cung bậc: từ lo âu, cạn kiệt năng lượng
cho đến khoảnh khắc vỡ òa hân hoan khi sản phẩm thực sự hoạt động.

Những câu chuyện hậu trường ``dở khóc dở cười'' như: tranh cãi nội bộ vì
chưa phân định rõ vai trò, thức trắng đêm debug code đến 3 giờ sáng,
uống cạn 5 lon bò húc, quên commit code hay vô tình push nhầm file chứa
mã bí mật (\texttt{.env}) lên GitHub. Tất cả đã để lại những bài học
``nhớ đời'' mà không một trường lớp hay sách vở nào có thể dạy được.

\hypertarget{buxe0i-hux1ecdc-ux111uxfac-kux1ebft-key-takeaways}{%
\subsection{4. Bài Học Đúc Kết (Key
Takeaways)}}

\hypertarget{tux1b0-duy-phuxe1t-triux1ec3n-sux1ea3n-phux1ea9m-tinh-gux1ecdn-pragmatic-mvp-mindset}{%
\subsubsection{Tư duy phát triển sản phẩm tinh gọn (Pragmatic MVP
Mindset)}}

``Small, finished work beats big, broken ideas'' (Một sản phẩm nhỏ nhưng
hoàn thiện luôn chiến thắng một ý tưởng vĩ đại nhưng đổ vỡ). Đây là bài
học đắt giá nhất để vượt qua áp lực thời gian: dũng cảm cắt gọt tính
năng và tập trung làm thật tốt một tính năng cốt lõi (Scope it tiny).

Để có một bài thuyết trình 3 phút thành công trước ban giám khảo, việc
phân chia vai trò rõ ràng (ai code, ai thiết kế, ai pitch) và chuẩn bị
kỹ lưỡng kịch bản demo chính là yếu tố mang tính quyết định.

\hypertarget{kiux1ebfn-truxfac-linh-houx1ea1t-vuxe0-tux1ed1i-ux1b0u-huxf3a-chi-phuxed-vux1eadn-huxe0nh-flexible-cost-efficient-architecture}{%
\subsubsection{Kiến trúc linh hoạt và tối ưu hóa chi phí vận hành
(Flexible \& Cost-Efficient
Architecture)}}

Thiết kế hệ thống thông minh là thiết kế cho phép sản phẩm thay đổi, mở
rộng mà không phải ``đập đi xây lại'' toàn bộ. Điều này đạt được nhờ
việc tận dụng các adapter kết nối kênh giao tiếp và phân tách các công
cụ xử lý logic độc lập.

Việc ứng dụng Framework như AgentCore không chỉ giúp cắt giảm đến 60\%
mã nguồn hạ tầng mà còn duy trì chi phí vận hành cực kỳ kinh tế (chỉ
\$0.006 cho mỗi đơn đặt hàng thông qua KFC Bot, hoặc khả năng co giãn
chi phí linh hoạt dựa trên tài nguyên Serverless như dự án Signal
Scout).

\hypertarget{giuxe1-trux1ecb-cux1ee7a-cux1ed9ng-ux111ux1ed3ng-vuxe0-sux1ef1-bux1ed5-khuyux1ebft-ux111ux1ed9i-nhuxf3m-the-power-of-people-diversity}{%
\subsubsection{Giá trị của cộng đồng và sự bổ khuyết đội nhóm (The Power
of People \&
Diversity)}}

Chiến thắng lớn nhất sau một kỳ Hackathon không nằm ở giải thưởng vật
chất, mà nằm ở những người đồng đội đã kề vai sát cánh chiến đấu trong
đêm, cùng những mối quan hệ bền chặt được kết nối trong cộng đồng.

Sự đa dạng về thế mạnh cá nhân trong một đội nhóm (khác biệt kỹ năng
luôn tốt hơn sự đồng đều kỹ năng) chính là chìa khóa để giải quyết bài
toán đa chiều dưới áp lực thời gian khủng khiếp.

\hypertarget{ux1ee9ng-dux1ee5ng-vuxe0o-thux1ef1c-tux1ebf-applying-to-work}{%
\subsection{5. Ứng Dụng Vào Thực Tế (Applying to
Work)}}

\begin{itemize}
\tightlist
\item
  \textbf{Áp dụng triết lý ``Scope it tiny''}: Luôn ưu tiên hoàn thành
  dứt điểm các phiên bản sản phẩm nhỏ (MVP) để lấy phản hồi thực tế từ
  người dùng, tránh sa đà vào việc thiết kế những tính năng đồ sộ nhưng
  không bao giờ hoàn thiện.
\item
  \textbf{Xây dựng kiến trúc hướng tác tử (Agent-based)}: Thiết kế hệ
  thống theo dạng module hóa, giúp dễ dàng cắm rút (plug-and-play) thêm
  các công cụ tự động hóa hoặc kênh phân phối mới mà không làm xáo trộn
  mã nguồn cốt lõi.
\item
  \textbf{Quản trị chi phí từ trong trứng nước}: Nghiên cứu sâu phương
  pháp tối ưu hóa chi phí điện toán đám mây dựa trên việc đo lường tài
  nguyên thực tế. Tích cực sử dụng các công cụ tính toán chi phí (AWS
  Pricing API/Calculator) ngay từ bước phác thảo thiết kế.
\item
  \textbf{Chuẩn bị sẵn sàng bệ phóng}: Luôn chủ động chuẩn bị sẵn các bộ
  khung dự án (starter templates), tài khoản Cloud và công cụ cần thiết
  để có thể đẩy nhanh tốc độ triển khai ngay khi có dự án mới phát sinh
  đột xuất.
\end{itemize}

\hypertarget{trux1ea3i-nghiux1ec7m-khuxf3a-hux1ecdc-event-experience}{%
\subsection{6. Trải Nghiệm Khóa Học (Event
Experience)}}

\begin{itemize}
\tightlist
\item
  \textbf{Trải nghiệm sống động từ diễn giả}: Những chia sẻ chân thực,
  mộc mạc và không hề giấu giếm về những thất bại hay sự hỗn loạn của
  các đội thi đã thắp sáng ngọn lửa tự tin, giúp tôi rũ bỏ nỗi sợ
  ``không đủ năng lực'' để mạnh dạn dấn thân.
\item
  \textbf{Góc nhìn công nghệ thực tế}: Việc tận mắt chứng kiến các mô
  hình như Amazon Bedrock, AgentCore phối hợp cùng Computer Vision
  (YOLO) để tạo nên những hệ thống tự vận hành xuất sắc đã giúp tôi định
  hình rõ nét xu hướng công nghệ của tương lai.
\item
  \textbf{Sử dụng công cụ thông minh}: Tôi hiểu sâu sắc hơn giá trị to
  lớn của AI Agent trong việc tự động hóa các tác vụ lặp đi lặp lại,
  nâng cao năng suất cá nhân lên gấp nhiều lần mà vẫn đảm bảo tính chính
  xác cao.
\item
  \textbf{Kết nối cộng đồng}: Không khí rực lửa, hào hứng và tinh thần
  hỗ trợ không biên giới của cộng đồng AWS tại sự kiện đã truyền cho tôi
  một nguồn năng lượng mạnh mẽ để tiếp tục học hỏi và cống hiến.
\end{itemize}

\hypertarget{tux1ed5ng-kux1ebft-lessons-learned}{%
\subsection{7. Tổng Kết (Lessons
Learned)}}

\begin{itemize}
\tightlist
\item
  Đừng bao giờ chờ đợi đến khi bản thân cảm thấy ``hoàn toàn sẵn sàng''
  mới dám đăng ký tham gia các sân chơi thực chiến. Sự dũng cảm bước ra
  khỏi vùng an toàn mới chính là bệ phóng lớn nhất.
\item
  Một hệ thống AI Agent xuất sắc không chỉ cần một mô hình ngôn ngữ lớn
  (LLM) thông minh, mà quan trọng hơn là khả năng kết nối chính xác với
  dữ liệu nội bộ đáng tin cậy của doanh nghiệp để đưa ra hành động thực
  tế.
\item
  Hãy học cách yêu thích sự hỗn loạn, các lỗi debug lúc nửa đêm và những
  sản phẩm ``chưa hoàn hảo''. Bởi vì đó chính là nơi lưu giữ những bài
  học sâu sắc nhất và những kỷ niệm vô giá nhất trong cuộc đời làm kỹ
  thuật.
\end{itemize}

\emph{(Hình ảnh sự kiện có thể được bổ sung tại đây)}

\begin{figure}
\centering
\includegraphics{event3.jpg}
\caption{Event 3}
\end{figure}
